%%%%%%%%%%%%%%%%%%%%%%%%%%%%%%%%%%%%%%%%%%%%%%%%%%%%%%%%%%%%%%%%%%%%%%%%%%%%%%%%
% CHAPTER 11
%%%%%%%%%%%%%%%%%%%%%%%%%%%%%%%%%%%%%%%%%%%%%%%%%%%%%%%%%%%%%%%%%%%%%%%%%%%%%%%%

\chapter{z-Transform}

\begin{tcolorbox}[colback=blue!5!white,colframe=blue!70!black,title=Learning Objectives]

After studying this chapter, the reader will be able to

\begin{itemize}
\item Understand the need for the z-transform.
\item Define the bilateral and unilateral z-transforms.
\item Determine the Region of Convergence (ROC).
\item Compute z-transforms of standard sequences.
\item Understand poles and zeros.
\item Relate the z-transform to the DTFT.
\end{itemize}

\end{tcolorbox}

%%%%%%%%%%%%%%%%%%%%%%%%%%%%%%%%%%%%%%%%%%%%%%%%%%%%%%%%%%%%%%%%%%%%%%%%%%%%%%%%

\section{Introduction}

The DTFT is useful for frequency analysis, but it does not exist for many important
sequences.

Example:

\[
x[n]=u[n]
\]

has a DTFT that is not absolutely convergent.

To analyze such sequences and solve difference equations systematically, we use the
\textbf{z-transform}.

The z-transform is the discrete-time counterpart of the Laplace transform.

%%%%%%%%%%%%%%%%%%%%%%%%%%%%%%%%%%%%%%%%%%%%%%%%%%%%%%%%%%%%%%%%%%%%%%%%%%%%%%%%

\section{Complex Variable \(z\)}

The complex variable is written as

\[
\boxed{
z=re^{j\omega}
}
\]

where

\begin{itemize}
\item \(r\): magnitude,
\item \(\omega\): angle.
\end{itemize}

The \(z\)-plane is a complex plane with

\[
z=\sigma+j\Omega.
\]

The unit circle is

\[
\boxed{
|z|=1.
}
\]

%%%%%%%%%%%%%%%%%%%%%%%%%%%%%%%%%%%%%%%%%%%%%%%%%%%%%%%%%%%%%%%%%%%%%%%%%%%%%%%%

\section{Bilateral z-Transform}

The bilateral (two-sided) z-transform is

\[
\boxed{
X(z)=
\sum_{n=-\infty}^{\infty}
x[n]z^{-n}
}
\]

The summation includes all integer values of \(n\).

%%%%%%%%%%%%%%%%%%%%%%%%%%%%%%%%%%%%%%%%%%%%%%%%%%%%%%%%%%%%%%%%%%%%%%%%%%%%%%%%

\section{Region of Convergence (ROC)}

The values of \(z\) for which the series converges form the
\textbf{Region of Convergence}.

\[
\boxed{
\text{ROC}=
\{z:\sum |x[n]z^{-n}|<\infty\}
}
\]

The ROC is as important as the transform itself.

%%%%%%%%%%%%%%%%%%%%%%%%%%%%%%%%%%%%%%%%%%%%%%%%%%%%%%%%%%%%%%%%%%%%%%%%%%%%%%%%

\section{Unilateral z-Transform}

For causal sequences, we often use the one-sided transform:

\[
\boxed{
X^+(z)=
\sum_{n=0}^{\infty}
x[n]z^{-n}
}
\]

This is convenient for solving difference equations with initial conditions.

%%%%%%%%%%%%%%%%%%%%%%%%%%%%%%%%%%%%%%%%%%%%%%%%%%%%%%%%%%%%%%%%%%%%%%%%%%%%%%%%

\section{z-Transform of the Unit Sample}

Let

\[
x[n]=\delta[n].
\]

Then

\[
X(z)=
\sum_{n=-\infty}^{\infty}
\delta[n]z^{-n}.
\]

Only \(n=0\) contributes:

\[
\boxed{
X(z)=1
}
\]

The ROC is the entire \(z\)-plane.

%%%%%%%%%%%%%%%%%%%%%%%%%%%%%%%%%%%%%%%%%%%%%%%%%%%%%%%%%%%%%%%%%%%%%%%%%%%%%%%%

\section{z-Transform of the Unit Step}

Let

\[
x[n]=u[n].
\]

Then

\[
X(z)=
\sum_{n=0}^{\infty}z^{-n}.
\]

This is a geometric series:

\[
X(z)=
\frac{1}{1-z^{-1}}
\]

provided

\[
|z^{-1}|<1.
\]

Hence,

\[
\boxed{
u[n]
\xleftrightarrow{\mathcal{Z}}
\frac{1}{1-z^{-1}},
\qquad
\text{ROC: }|z|>1
}
\]

%%%%%%%%%%%%%%%%%%%%%%%%%%%%%%%%%%%%%%%%%%%%%%%%%%%%%%%%%%%%%%%%%%%%%%%%%%%%%%%%

\section{z-Transform of \(a^n u[n]\)}

Let

\[
x[n]=a^n u[n].
\]

Then

\[
X(z)=
\sum_{n=0}^{\infty}
a^n z^{-n}
=
\sum_{n=0}^{\infty}
(az^{-1})^n.
\]

Therefore,

\[
\boxed{
X(z)=
\frac{1}{1-az^{-1}},
\qquad
\text{ROC: }|z|>|a|
}
\]

%%%%%%%%%%%%%%%%%%%%%%%%%%%%%%%%%%%%%%%%%%%%%%%%%%%%%%%%%%%%%%%%%%%%%%%%%%%%%%%%

\section{Example}

For

\[
x[n]=\left(\frac12\right)^n u[n],
\]

\[
\boxed{
X(z)=
\frac{1}{1-\frac12 z^{-1}},
\qquad
\text{ROC: }|z|>\frac12
}
\]

%%%%%%%%%%%%%%%%%%%%%%%%%%%%%%%%%%%%%%%%%%%%%%%%%%%%%%%%%%%%%%%%%%%%%%%%%%%%%%%%

\section{Left-Sided Sequence}

Consider

\[
x[n]=-a^n u[-n-1].
\]

Then

\[
X(z)=
\sum_{n=-\infty}^{-1}
(-a^n)z^{-n}.
\]

After simplification,

\[
\boxed{
X(z)=
\frac{1}{1-az^{-1}},
\qquad
\text{ROC: }|z|<|a|
}
\]

Notice that the algebraic expression is the same as the right-sided sequence, but
the ROC is different.

%%%%%%%%%%%%%%%%%%%%%%%%%%%%%%%%%%%%%%%%%%%%%%%%%%%%%%%%%%%%%%%%%%%%%%%%%%%%%%%%

\section{ROC Determines the Sequence}

\[
\frac{1}{1-az^{-1}}
\]

can represent three different sequences:

\begin{table}[H]
\centering
\caption{Same transform, different ROCs}
\begin{tabular}{|p{5cm}|p{5cm}|}
\hline
\textbf{Sequence} & \textbf{ROC}\\
\hline
\(a^n u[n]\) & \(|z|>|a|\)\\
\hline
\(-a^n u[-n-1]\) & \(|z|<|a|\)\\
\hline
No two-sided sequence exists & \(|z|=|a|\)\\
\hline
\end{tabular}
\end{table}

%%%%%%%%%%%%%%%%%%%%%%%%%%%%%%%%%%%%%%%%%%%%%%%%%%%%%%%%%%%%%%%%%%%%%%%%%%%%%%%%

\section{Poles and Zeros}

%%%%%%%%%%%%%%%%%%%%%%%%%%%%%%%%%%%%%%%%%%%%%%%%%%%%%%%%%%%%%%%%%%%%%%%%%%%%%%%%

\subsection{Zero}

A value of \(z\) that makes

\[
X(z)=0.
\]

%%%%%%%%%%%%%%%%%%%%%%%%%%%%%%%%%%%%%%%%%%%%%%%%%%%%%%%%%%%%%%%%%%%%%%%%%%%%%%%%

\subsection{Pole}

A value of \(z\) that makes

\[
X(z)\to\infty.
\]

%%%%%%%%%%%%%%%%%%%%%%%%%%%%%%%%%%%%%%%%%%%%%%%%%%%%%%%%%%%%%%%%%%%%%%%%%%%%%%%%

\subsection{Example}

\[
X(z)=
\frac{1-z^{-1}}
{1-0.5z^{-1}}.
\]

Zero:

\[
z=1.
\]

Pole:

\[
z=0.5.
\]

%%%%%%%%%%%%%%%%%%%%%%%%%%%%%%%%%%%%%%%%%%%%%%%%%%%%%%%%%%%%%%%%%%%%%%%%%%%%%%%%

\section{Pole-Zero Plot}

\begin{center}
\begin{tikzpicture}[scale=2]
\draw[->] (-1.5,0) -- (1.5,0) node[right] {Re};
\draw[->] (0,-1.2) -- (0,1.2) node[above] {Im};

\draw (0,0) circle (1);

\draw[thick] (1,0) node[circle,draw,inner sep=1.5pt] {};
\draw[thick] (0.5,0) node[cross out,draw,inner sep=2pt] {};

\node[below] at (1,0) {$z=1$};
\node[below] at (0.5,0) {$z=0.5$};
\end{tikzpicture}
\end{center}

Convention:

\begin{itemize}
\item \(\circ\): zero,
\item \(\times\): pole.
\end{itemize}

%%%%%%%%%%%%%%%%%%%%%%%%%%%%%%%%%%%%%%%%%%%%%%%%%%%%%%%%%%%%%%%%%%%%%%%%%%%%%%%%

\section{ROC of Rational Functions}

For rational functions, the ROC is an annular region.

%%%%%%%%%%%%%%%%%%%%%%%%%%%%%%%%%%%%%%%%%%%%%%%%%%%%%%%%%%%%%%%%%%%%%%%%%%%%%%%%

\subsection{Example}

\[
X(z)=
\frac{1}
{(1-0.5z^{-1})(1-2z^{-1})}.
\]

Possible ROCs:

\begin{itemize}
\item \(|z|>2\),
\item \(0.5<|z|<2\),
\item \(|z|<0.5\).
\end{itemize}

The ROC cannot contain poles.

%%%%%%%%%%%%%%%%%%%%%%%%%%%%%%%%%%%%%%%%%%%%%%%%%%%%%%%%%%%%%%%%%%%%%%%%%%%%%%%%

\section{Causality from ROC}

A rational sequence is causal if

\[
\boxed{
\text{ROC is outside the outermost pole.}
}
\]

Example:

\[
X(z)=
\frac{1}{1-0.5z^{-1}},
\qquad
\text{ROC: }|z|>0.5.
\]

This corresponds to

\[
x[n]=(0.5)^n u[n]
\]

which is causal.

%%%%%%%%%%%%%%%%%%%%%%%%%%%%%%%%%%%%%%%%%%%%%%%%%%%%%%%%%%%%%%%%%%%%%%%%%%%%%%%%

\section{Stability from ROC}

A sequence is stable if the DTFT exists, which requires the ROC to include the unit circle.

\[
\boxed{
|z|=1 \in \text{ROC}
}
\]

For the previous example,

\[
|z|>0.5
\]

includes \(|z|=1\), so the sequence is stable.

%%%%%%%%%%%%%%%%%%%%%%%%%%%%%%%%%%%%%%%%%%%%%%%%%%%%%%%%%%%%%%%%%%%%%%%%%%%%%%%%

\section{Connection with DTFT}

Evaluate the z-transform on the unit circle:

\[
\boxed{
X(e^{j\omega})=
X(z)\Big|_{z=e^{j\omega}}
}
\]

This is valid only if the unit circle lies inside the ROC.

%%%%%%%%%%%%%%%%%%%%%%%%%%%%%%%%%%%%%%%%%%%%%%%%%%%%%%%%%%%%%%%%%%%%%%%%%%%%%%%%

\section{Example}

\[
X(z)=
\frac{1}{1-0.5z^{-1}},
\qquad
|z|>0.5.
\]

Since the unit circle is included,

\[
X(e^{j\omega})=
\frac{1}{1-0.5e^{-j\omega}}.
\]

This is exactly the DTFT of \((0.5)^n u[n]\).

%%%%%%%%%%%%%%%%%%%%%%%%%%%%%%%%%%%%%%%%%%%%%%%%%%%%%%%%%%%%%%%%%%%%%%%%%%%%%%%%

\section{Initial Value Theorem}

For causal sequences,

\[
\boxed{
x[0]=
\lim_{z\to\infty}X(z)
}
\]

%%%%%%%%%%%%%%%%%%%%%%%%%%%%%%%%%%%%%%%%%%%%%%%%%%%%%%%%%%%%%%%%%%%%%%%%%%%%%%%%

\subsection{Example}

\[
X(z)=
\frac{1}{1-0.5z^{-1}}.
\]

Then

\[
x[0]=
\lim_{z\to\infty}
\frac{1}{1-0.5z^{-1}}
=1.
\]

%%%%%%%%%%%%%%%%%%%%%%%%%%%%%%%%%%%%%%%%%%%%%%%%%%%%%%%%%%%%%%%%%%%%%%%%%%%%%%%%

\section{Final Value Theorem}

If the limits exist and all poles of \((1-z^{-1})X(z)\) are inside the unit circle,

\[
\boxed{
\lim_{n\to\infty}x[n]
=
\lim_{z\to1}(1-z^{-1})X(z)
}
\]

%%%%%%%%%%%%%%%%%%%%%%%%%%%%%%%%%%%%%%%%%%%%%%%%%%%%%%%%%%%%%%%%%%%%%%%%%%%%%%%%

\subsection{Example}

\[
X(z)=
\frac{1}{1-0.5z^{-1}}.
\]

\[
(1-z^{-1})X(z)=
\frac{1-z^{-1}}{1-0.5z^{-1}}.
\]

At \(z=1\),

\[
\lim_{n\to\infty}x[n]=0.
\]

This matches

\[
(0.5)^n\to0.
\]

%%%%%%%%%%%%%%%%%%%%%%%%%%%%%%%%%%%%%%%%%%%%%%%%%%%%%%%%%%%%%%%%%%%%%%%%%%%%%%%%

\section{Common ROC Patterns}

\begin{table}[H]
\centering
\caption{ROC patterns}
\begin{tabular}{|p{4cm}|p{5cm}|}
\hline
\textbf{Sequence Type} & \textbf{ROC}\\
\hline
Right-sided & Outside outermost pole\\
\hline
Left-sided & Inside innermost pole\\
\hline
Two-sided & Between poles\\
\hline
Finite-duration right-sided & Entire plane except \(z=0\)\\
\hline
Finite-duration left-sided & Entire plane except \(z=\infty\)\\
\hline
\end{tabular}
\end{table}

%%%%%%%%%%%%%%%%%%%%%%%%%%%%%%%%%%%%%%%%%%%%%%%%%%%%%%%%%%%%%%%%%%%%%%%%%%%%%%%%

\section{Solved Example 11.1}

Find the z-transform of

\[
x[n]=3^n u[n].
\]

%%%%%%%%%%%%%%%%%%%%%%%%%%%%%%%%%%%%%%%%%%%%%%%%%%%%%%%%%%%%%%%%%%%%%%%%%%%%%%%%

\subsection{Solution}

\[
X(z)=
\sum_{n=0}^{\infty}(3z^{-1})^n
=
\frac{1}{1-3z^{-1}}
\]

with ROC

\[
\boxed{
|z|>3
}
\]

%%%%%%%%%%%%%%%%%%%%%%%%%%%%%%%%%%%%%%%%%%%%%%%%%%%%%%%%%%%%%%%%%%%%%%%%%%%%%%%%

\section{Solved Example 11.2}

Find the z-transform of

\[
x[n]=-\left(\frac12\right)^n u[-n-1].
\]

%%%%%%%%%%%%%%%%%%%%%%%%%%%%%%%%%%%%%%%%%%%%%%%%%%%%%%%%%%%%%%%%%%%%%%%%%%%%%%%%

\subsection{Solution}

\[
\boxed{
X(z)=
\frac{1}{1-\frac12 z^{-1}},
\qquad
\text{ROC: }|z|<\frac12
}
\]

%%%%%%%%%%%%%%%%%%%%%%%%%%%%%%%%%%%%%%%%%%%%%%%%%%%%%%%%%%%%%%%%%%%%%%%%%%%%%%%%

\section{Solved Example 11.3}

Determine whether the sequence is stable:

\[
X(z)=
\frac{1}{1-2z^{-1}},
\qquad
\text{ROC: }|z|>2.
\]

%%%%%%%%%%%%%%%%%%%%%%%%%%%%%%%%%%%%%%%%%%%%%%%%%%%%%%%%%%%%%%%%%%%%%%%%%%%%%%%%

\subsection{Solution}

The unit circle \(|z|=1\) is not included in the ROC.

Hence,

\[
\boxed{
\text{The sequence is unstable.}
}
\]

%%%%%%%%%%%%%%%%%%%%%%%%%%%%%%%%%%%%%%%%%%%%%%%%%%%%%%%%%%%%%%%%%%%%%%%%%%%%%%%%

\section{Solved Example 11.4}

Find the poles and zeros of

\[
X(z)=
\frac{1-z^{-1}}
{1-0.25z^{-1}}.
\]

%%%%%%%%%%%%%%%%%%%%%%%%%%%%%%%%%%%%%%%%%%%%%%%%%%%%%%%%%%%%%%%%%%%%%%%%%%%%%%%%

\subsection{Solution}

Zero:

\[
\boxed{
z=1
}
\]

Pole:

\[
\boxed{
z=0.25
}
\]

%%%%%%%%%%%%%%%%%%%%%%%%%%%%%%%%%%%%%%%%%%%%%%%%%%%%%%%%%%%%%%%%%%%%%%%%%%%%%%%%

\section{Quick Check}

\begin{enumerate}
\item \(u[n]\rightarrow \boxed{\dfrac{1}{1-z^{-1}}}\).
\item ROC of \(u[n]\): \(\boxed{|z|>1}\).
\item Pole of \(\dfrac{1}{1-0.5z^{-1}}\): \(\boxed{0.5}\).
\item Stable if: \(\boxed{\text{unit circle is in ROC}}\).
\item DTFT is obtained by: \(\boxed{z=e^{j\omega}}\).
\end{enumerate}

%%%%%%%%%%%%%%%%%%%%%%%%%%%%%%%%%%%%%%%%%%%%%%%%%%%%%%%%%%%%%%%%%%%%%%%%%%%%%%%%

\begin{tcolorbox}[title=One-Minute Revision,colback=blue!5]

\[
X(z)=
\sum_{n=-\infty}^{\infty}x[n]z^{-n}
\]

\[
X^+(z)=
\sum_{n=0}^{\infty}x[n]z^{-n}
\]

\[
u[n]\leftrightarrow
\frac{1}{1-z^{-1}},
\quad |z|>1
\]

\[
a^n u[n]\leftrightarrow
\frac{1}{1-az^{-1}},
\quad |z|>|a|
\]

Pole: denominator zero.

Zero: numerator zero.

Causal:

\[
\text{ROC outside outermost pole}
\]

Stable:

\[
|z|=1 \in \text{ROC}
\]

DTFT:

\[
X(e^{j\omega})=X(z)|_{z=e^{j\omega}}
\]

\end{tcolorbox}
%%%%%%%%%%%%%%%%%%%%%%%%%%%%%%%%%%%%%%%%%%%%%%%%%%%%%%%%%%%%%%%%%%%%%%%%%%%%%%%%
%%%%%%%%%%%%%%%%%%%%%%%%%%%%%%%%%%%%%%%%%%%%%%%%%%%%%%%%%%%%%%%%%%%%%%%%%%%%%%%%
\section{Linearity Property}

If

\[
x_1[n]\xleftrightarrow{\mathcal{Z}}X_1(z)
\]

and

\[
x_2[n]\xleftrightarrow{\mathcal{Z}}X_2(z),
\]

then

\[
\boxed{
a_1x_1[n]+a_2x_2[n]
\xleftrightarrow{\mathcal{Z}}
a_1X_1(z)+a_2X_2(z)
}
\]

%%%%%%%%%%%%%%%%%%%%%%%%%%%%%%%%%%%%%%%%%%%%%%%%%%%%%%%%%%%%%%%%%%%%%%%%%%%%%%%%

\section{Time-Shifting Property}

%%%%%%%%%%%%%%%%%%%%%%%%%%%%%%%%%%%%%%%%%%%%%%%%%%%%%%%%%%%%%%%%%%%%%%%%%%%%%%%%

\subsection{Delay}

\[
\boxed{
x[n-n_0]
\xleftrightarrow{\mathcal{Z}}
z^{-n_0}X(z)
}
\]

For bilateral transforms, the ROC is unchanged except possibly at the origin or infinity.

%%%%%%%%%%%%%%%%%%%%%%%%%%%%%%%%%%%%%%%%%%%%%%%%%%%%%%%%%%%%%%%%%%%%%%%%%%%%%%%%

\subsection{Advance}

\[
\boxed{
x[n+n_0]
\xleftrightarrow{\mathcal{Z}}
z^{n_0}X(z)
}
\]

with additional terms for unilateral transforms.

%%%%%%%%%%%%%%%%%%%%%%%%%%%%%%%%%%%%%%%%%%%%%%%%%%%%%%%%%%%%%%%%%%%%%%%%%%%%%%%%

\section{Scaling Property}

\[
\boxed{
a^n x[n]
\xleftrightarrow{\mathcal{Z}}
X\left(\frac{z}{a}\right)
}
\]

%%%%%%%%%%%%%%%%%%%%%%%%%%%%%%%%%%%%%%%%%%%%%%%%%%%%%%%%%%%%%%%%%%%%%%%%%%%%%%%%

\section{Time Reversal}

\[
\boxed{
x[-n]
\xleftrightarrow{\mathcal{Z}}
X(z^{-1})
}
\]

The ROC is inverted with respect to the unit circle.

%%%%%%%%%%%%%%%%%%%%%%%%%%%%%%%%%%%%%%%%%%%%%%%%%%%%%%%%%%%%%%%%%%%%%%%%%%%%%%%%

\section{Convolution Property}

If

\[
y[n]=x[n]*h[n],
\]

then

\[
\boxed{
Y(z)=X(z)H(z)
}
\]

This is one of the most powerful properties of the z-transform.

%%%%%%%%%%%%%%%%%%%%%%%%%%%%%%%%%%%%%%%%%%%%%%%%%%%%%%%%%%%%%%%%%%%%%%%%%%%%%%%%

\section{Differentiation Property}

\[
\boxed{
nx[n]
\xleftrightarrow{\mathcal{Z}}
-z\frac{dX(z)}{dz}
}
\]

This is useful for finding moments of sequences.

%%%%%%%%%%%%%%%%%%%%%%%%%%%%%%%%%%%%%%%%%%%%%%%%%%%%%%%%%%%%%%%%%%%%%%%%%%%%%%%%

\section{Initial Value Theorem}

For causal sequences,

\[
\boxed{
x[0]=\lim_{z\to\infty}X(z)
}
\]

%%%%%%%%%%%%%%%%%%%%%%%%%%%%%%%%%%%%%%%%%%%%%%%%%%%%%%%%%%%%%%%%%%%%%%%%%%%%%%%%

\subsection{Example}

\[
X(z)=
\frac{1}{1-0.5z^{-1}}.
\]

Then

\[
x[0]=
\lim_{z\to\infty}
\frac{1}{1-0.5z^{-1}}
=1.
\]

%%%%%%%%%%%%%%%%%%%%%%%%%%%%%%%%%%%%%%%%%%%%%%%%%%%%%%%%%%%%%%%%%%%%%%%%%%%%%%%%

\section{Final Value Theorem}

If all poles of \((1-z^{-1})X(z)\) are inside the unit circle,

\[
\boxed{
\lim_{n\to\infty}x[n]
=
\lim_{z\to1}(1-z^{-1})X(z)
}
\]

%%%%%%%%%%%%%%%%%%%%%%%%%%%%%%%%%%%%%%%%%%%%%%%%%%%%%%%%%%%%%%%%%%%%%%%%%%%%%%%%

\subsection{Example}

\[
X(z)=
\frac{1}{1-0.5z^{-1}}.
\]

\[
(1-z^{-1})X(z)=
\frac{1-z^{-1}}{1-0.5z^{-1}}.
\]

At \(z=1\),

\[
\lim_{n\to\infty}x[n]=0.
\]

%%%%%%%%%%%%%%%%%%%%%%%%%%%%%%%%%%%%%%%%%%%%%%%%%%%%%%%%%%%%%%%%%%%%%%%%%%%%%%%%

\section{Inverse z-Transform}

We now recover \(x[n]\) from \(X(z)\).

%%%%%%%%%%%%%%%%%%%%%%%%%%%%%%%%%%%%%%%%%%%%%%%%%%%%%%%%%%%%%%%%%%%%%%%%%%%%%%%%

\subsection{Method 1: Power-Series Expansion}

Given

\[
X(z)=
\frac{1}{1-0.5z^{-1}},
\qquad |z|>0.5,
\]

expand as

\[
X(z)=
1+0.5z^{-1}+0.25z^{-2}+\cdots
\]

Hence,

\[
\boxed{
x[n]=(0.5)^n u[n]
}
\]

%%%%%%%%%%%%%%%%%%%%%%%%%%%%%%%%%%%%%%%%%%%%%%%%%%%%%%%%%%%%%%%%%%%%%%%%%%%%%%%%

\section{Method 2: Partial-Fraction Expansion}

Consider

\[
X(z)=
\frac{1}
{(1-0.5z^{-1})(1-0.25z^{-1})}.
\]

Assume

\[
X(z)=
\frac{A}{1-0.5z^{-1}}
+
\frac{B}{1-0.25z^{-1}}.
\]

%%%%%%%%%%%%%%%%%%%%%%%%%%%%%%%%%%%%%%%%%%%%%%%%%%%%%%%%%%%%%%%%%%%%%%%%%%%%%%%%

\subsection{Find \(A\) and \(B\)}

Multiplying both sides,

\[
1=A(1-0.25z^{-1})+B(1-0.5z^{-1}).
\]

Set

\[
z^{-1}=2
\]

to eliminate the first denominator:

\[
1=B(1-1)=0
\]

Instead, compare coefficients:

\[
A+B=1
\]

\[
0.25A+0.5B=0.
\]

Solving,

\[
A=2,
\qquad
B=-1.
\]

%%%%%%%%%%%%%%%%%%%%%%%%%%%%%%%%%%%%%%%%%%%%%%%%%%%%%%%%%%%%%%%%%%%%%%%%%%%%%%%%

\subsection{Inverse Transform}

\[
\boxed{
x[n]=
2(0.5)^n u[n]
-(0.25)^n u[n]
}
\]

%%%%%%%%%%%%%%%%%%%%%%%%%%%%%%%%%%%%%%%%%%%%%%%%%%%%%%%%%%%%%%%%%%%%%%%%%%%%%%%%

\section{Repeated Poles}

If

\[
X(z)=
\frac{1}{(1-az^{-1})^2},
\]

then

\[
\boxed{
x[n]=(n+1)a^n u[n]
}
\]

Repeated poles produce polynomial factors in time.

%%%%%%%%%%%%%%%%%%%%%%%%%%%%%%%%%%%%%%%%%%%%%%%%%%%%%%%%%%%%%%%%%%%%%%%%%%%%%%%%

\section{Difference Equations}

Consider the first-order equation

\[
y[n]-0.5y[n-1]=x[n].
\]

%%%%%%%%%%%%%%%%%%%%%%%%%%%%%%%%%%%%%%%%%%%%%%%%%%%%%%%%%%%%%%%%%%%%%%%%%%%%%%%%

\subsection{Zero Initial Conditions}

Taking the unilateral z-transform,

\[
Y(z)-0.5z^{-1}Y(z)=X(z).
\]

Hence,

\[
Y(z)=
\frac{X(z)}{1-0.5z^{-1}}.
\]

Therefore,

\[
\boxed{
H(z)=
\frac{1}{1-0.5z^{-1}}
}
\]

%%%%%%%%%%%%%%%%%%%%%%%%%%%%%%%%%%%%%%%%%%%%%%%%%%%%%%%%%%%%%%%%%%%%%%%%%%%%%%%%

\section{Impulse Response}

For

\[
x[n]=\delta[n],
\]

\[
X(z)=1.
\]

Thus,

\[
H(z)=
\frac{1}{1-0.5z^{-1}}.
\]

Inverse transform:

\[
\boxed{
h[n]=(0.5)^n u[n]
}
\]

%%%%%%%%%%%%%%%%%%%%%%%%%%%%%%%%%%%%%%%%%%%%%%%%%%%%%%%%%%%%%%%%%%%%%%%%%%%%%%%%

\section{Step Response}

For

\[
x[n]=u[n],
\]

\[
X(z)=
\frac{1}{1-z^{-1}}.
\]

Then

\[
Y(z)=
\frac{1}{(1-z^{-1})(1-0.5z^{-1})}.
\]

Using partial fractions,

\[
Y(z)=
\frac{2}{1-z^{-1}}
-
\frac{1}{1-0.5z^{-1}}.
\]

Hence,

\[
\boxed{
y[n]=
2u[n]-(0.5)^n u[n]
}
\]

%%%%%%%%%%%%%%%%%%%%%%%%%%%%%%%%%%%%%%%%%%%%%%%%%%%%%%%%%%%%%%%%%%%%%%%%%%%%%%%%

\section{System Function}

For a general difference equation

\[
\sum_{k=0}^{N}a_k y[n-k]
=
\sum_{k=0}^{M}b_k x[n-k],
\]

taking the z-transform gives

\[
H(z)=
\frac{Y(z)}{X(z)}
=
\frac{\sum_{k=0}^{M}b_k z^{-k}}
{\sum_{k=0}^{N}a_k z^{-k}}.
\]

%%%%%%%%%%%%%%%%%%%%%%%%%%%%%%%%%%%%%%%%%%%%%%%%%%%%%%%%%%%%%%%%%%%%%%%%%%%%%%%%

\section{Pole-Zero Interpretation}

%%%%%%%%%%%%%%%%%%%%%%%%%%%%%%%%%%%%%%%%%%%%%%%%%%%%%%%%%%%%%%%%%%%%%%%%%%%%%%%%

\subsection{Poles Near the Unit Circle}

\begin{itemize}
\item Produce sharp resonances.
\item Give long impulse responses.
\item Make the system more selective.
\end{itemize}

%%%%%%%%%%%%%%%%%%%%%%%%%%%%%%%%%%%%%%%%%%%%%%%%%%%%%%%%%%%%%%%%%%%%%%%%%%%%%%%%

\subsection{Zeros on the Unit Circle}

\begin{itemize}
\item Create frequency nulls.
\item Completely reject specific frequencies.
\end{itemize}

%%%%%%%%%%%%%%%%%%%%%%%%%%%%%%%%%%%%%%%%%%%%%%%%%%%%%%%%%%%%%%%%%%%%%%%%%%%%%%%%

\section{Example: Notch Filter}

\[
H(z)=
\frac{1-2\cos\omega_0 z^{-1}+z^{-2}}
{1-r(2\cos\omega_0)z^{-1}+r^2 z^{-2}}.
\]

Zeros at

\[
z=e^{\pm j\omega_0}
\]

create a notch at frequency \(\omega_0\).

%%%%%%%%%%%%%%%%%%%%%%%%%%%%%%%%%%%%%%%%%%%%%%%%%%%%%%%%%%%%%%%%%%%%%%%%%%%%%%%%

\section{Stability Criterion}

A causal rational system is stable if all poles lie strictly inside the unit circle:

\[
\boxed{
|p_i|<1
}
\]

%%%%%%%%%%%%%%%%%%%%%%%%%%%%%%%%%%%%%%%%%%%%%%%%%%%%%%%%%%%%%%%%%%%%%%%%%%%%%%%%

\section{Example}

\[
H(z)=
\frac{1}{1-1.2z^{-1}}.
\]

Pole:

\[
z=1.2.
\]

Since

\[
|1.2|>1,
\]

\[
\boxed{
\text{The system is unstable.}
}
\]

%%%%%%%%%%%%%%%%%%%%%%%%%%%%%%%%%%%%%%%%%%%%%%%%%%%%%%%%%%%%%%%%%%%%%%%%%%%%%%%%

\section{Solved Example 11.5}

Find the inverse z-transform of

\[
X(z)=
\frac{1}{1-0.8z^{-1}},
\qquad |z|>0.8.
\]

%%%%%%%%%%%%%%%%%%%%%%%%%%%%%%%%%%%%%%%%%%%%%%%%%%%%%%%%%%%%%%%%%%%%%%%%%%%%%%%%

\subsection{Solution}

Using the standard pair,

\[
\boxed{
x[n]=(0.8)^n u[n]
}
\]

%%%%%%%%%%%%%%%%%%%%%%%%%%%%%%%%%%%%%%%%%%%%%%%%%%%%%%%%%%%%%%%%%%%%%%%%%%%%%%%%

\section{Solved Example 11.6}

Find the inverse z-transform of

\[
X(z)=
\frac{1}{1-0.8z^{-1}},
\qquad |z|<0.8.
\]

%%%%%%%%%%%%%%%%%%%%%%%%%%%%%%%%%%%%%%%%%%%%%%%%%%%%%%%%%%%%%%%%%%%%%%%%%%%%%%%%

\subsection{Solution}

This is a left-sided sequence:

\[
\boxed{
x[n]=-(0.8)^n u[-n-1]
}
\]

%%%%%%%%%%%%%%%%%%%%%%%%%%%%%%%%%%%%%%%%%%%%%%%%%%%%%%%%%%%%%%%%%%%%%%%%%%%%%%%%

\section{Solved Example 11.7}

Solve

\[
y[n]-0.5y[n-1]=u[n],
\qquad y[-1]=0.
\]

%%%%%%%%%%%%%%%%%%%%%%%%%%%%%%%%%%%%%%%%%%%%%%%%%%%%%%%%%%%%%%%%%%%%%%%%%%%%%%%%

\subsection{Step 1: Transform}

\[
Y(z)-0.5z^{-1}Y(z)=
\frac{1}{1-z^{-1}}.
\]

%%%%%%%%%%%%%%%%%%%%%%%%%%%%%%%%%%%%%%%%%%%%%%%%%%%%%%%%%%%%%%%%%%%%%%%%%%%%%%%%

\subsection{Step 2: Solve for \(Y(z)\)}

\[
Y(z)=
\frac{1}
{(1-z^{-1})(1-0.5z^{-1})}.
\]

%%%%%%%%%%%%%%%%%%%%%%%%%%%%%%%%%%%%%%%%%%%%%%%%%%%%%%%%%%%%%%%%%%%%%%%%%%%%%%%%

\subsection{Step 3: Partial Fractions}

\[
Y(z)=
\frac{2}{1-z^{-1}}
-
\frac{1}{1-0.5z^{-1}}.
\]

%%%%%%%%%%%%%%%%%%%%%%%%%%%%%%%%%%%%%%%%%%%%%%%%%%%%%%%%%%%%%%%%%%%%%%%%%%%%%%%%

\subsection{Step 4: Inverse Transform}

\[
\boxed{
y[n]=
2u[n]-(0.5)^n u[n]
}
\]

%%%%%%%%%%%%%%%%%%%%%%%%%%%%%%%%%%%%%%%%%%%%%%%%%%%%%%%%%%%%%%%%%%%%%%%%%%%%%%%%

\section{Solved Example 11.8}

Determine causality and stability:

\[
H(z)=
\frac{1}{(1-0.5z^{-1})(1-0.25z^{-1})}.
\]

%%%%%%%%%%%%%%%%%%%%%%%%%%%%%%%%%%%%%%%%%%%%%%%%%%%%%%%%%%%%%%%%%%%%%%%%%%%%%%%%

\subsection{Causality}

For a causal system,

\[
\text{ROC: }|z|>0.5.
\]

%%%%%%%%%%%%%%%%%%%%%%%%%%%%%%%%%%%%%%%%%%%%%%%%%%%%%%%%%%%%%%%%%%%%%%%%%%%%%%%%

\subsection{Stability}

The unit circle is included because

\[
1>0.5.
\]

Therefore,

\[
\boxed{
\text{The system is causal and stable.}
}
\]

%%%%%%%%%%%%%%%%%%%%%%%%%%%%%%%%%%%%%%%%%%%%%%%%%%%%%%%%%%%%%%%%%%%%%%%%%%%%%%%%

\section{Quick Comparison}

\begin{table}[H]
\centering
\caption{DTFT vs z-transform}
\begin{tabular}{|p{5cm}|p{4cm}|p{4cm}|}
\hline
\textbf{Feature} & \textbf{DTFT} & \textbf{z-Transform}\\
\hline
Variable & \(e^{j\omega}\) & \(z\)\\
\hline
Frequency analysis & Yes & Yes (via unit circle)\\
\hline
Difference equations & Limited & Excellent\\
\hline
ROC information & No & Yes\\
\hline
Stability analysis & Indirect & Direct\\
\hline
\end{tabular}
\end{table}

%%%%%%%%%%%%%%%%%%%%%%%%%%%%%%%%%%%%%%%%%%%%%%%%%%%%%%%%%%%%%%%%%%%%%%%%%%%%%%%%

\section{Important GATE Points}

\begin{enumerate}
\item The ROC is part of the complete answer.
\item Same algebraic expression can represent different sequences.
\item Causal \(\Rightarrow\) ROC outside outermost pole.
\item Stable \(\Rightarrow\) unit circle included in ROC.
\item Repeated poles produce polynomial terms.
\item Use unilateral z-transform when initial conditions are given.
\item DTFT is obtained by setting \(z=e^{j\omega}\).
\end{enumerate}

%%%%%%%%%%%%%%%%%%%%%%%%%%%%%%%%%%%%%%%%%%%%%%%%%%%%%%%%%%%%%%%%%%%%%%%%%%%%%%%%

\section{Chapter Summary}

In this chapter, we studied the properties of the z-transform, inverse z-transform
methods, initial and final value theorems, and the use of the unilateral
z-transform for solving difference equations.

We also interpreted poles and zeros geometrically and established the direct
connection between pole locations, causality and stability.

These ideas form the foundation for digital filter design, control systems and
advanced DSP analysis.

%%%%%%%%%%%%%%%%%%%%%%%%%%%%%%%%%%%%%%%%%%%%%%%%%%%%%%%%%%%%%%%%%%%%%%%%%%%%%%%%

\begin{tcolorbox}[title=One-Minute Revision,colback=blue!5]

\[
X(z)=
\sum_{n=-\infty}^{\infty}x[n]z^{-n}
\]

\[
x[n-n_0]\leftrightarrow z^{-n_0}X(z)
\]

\[
a^n x[n]\leftrightarrow X\left(\frac{z}{a}\right)
\]

\[
x*h\leftrightarrow X(z)H(z)
\]

\[
x[0]=\lim_{z\to\infty}X(z)
\]

\[
\lim_{n\to\infty}x[n]
=
\lim_{z\to1}(1-z^{-1})X(z)
\]

\[
a^n u[n]\leftrightarrow
\frac{1}{1-az^{-1}},
\quad |z|>|a|
\]

\[
(n+1)a^n u[n]
\leftrightarrow
\frac{1}{(1-az^{-1})^2}
\]

Causal:

\[
\text{ROC outside outermost pole}
\]

Stable:

\[
|p_i|<1
\]

DTFT:

\[
X(e^{j\omega})=X(z)|_{z=e^{j\omega}}
\]

\end{tcolorbox}
%%%%%%%%%%%%%%%%%%%%%%%%%%%%%%%%%%%%%%%%%%%%%%%%%%%%%%%%%%%%%%%%%%%%%%%%%%%%%%%%